\documentclass[10pt,a4paper]{amsart}
\usepackage{amsmath,amssymb,amsthm,graphicx}
% Retained minimal preamble for the Stacks Project chapter coherent.tex (Stacks chapter, pre-delegation retained fragment for stacks-coherent).
% Verbatim excerpt of the chapter preamble (sources/stacks-spaces-descent/preamble.tex): the reference/slogan/history
% environment block (native lines 7-12), Xy-pic setup (native lines 14-19) and the native theorem declarations with
% the mathematical macros (native lines 154-217). Copied mechanically; nothing added, removed or redefined. The
% wrapper class and the amsmath/amssymb/amsthm/graphicx base packages are supplied by the assembler, so the
% class/xr-hyper/hyperref/externaldocument part of the native preamble is omitted.
% For dealing with references we use the comment environment
\usepackage{verbatim}
\newenvironment{reference}{\comment}{\endcomment}
%\newenvironment{reference}{}{}
\newenvironment{slogan}{\comment}{\endcomment}
\newenvironment{history}{\comment}{\endcomment}
% For commutative diagrams we use Xy-pic
\usepackage[all]{xy}

% We use 2cell for 2-commutative diagrams.
\xyoption{2cell}
\UseAllTwocells
% Theorem environments.
%
\theoremstyle{plain}
\newtheorem{theorem}[subsection]{Theorem}
\newtheorem{proposition}[subsection]{Proposition}
\newtheorem{lemma}[subsection]{Lemma}

\theoremstyle{definition}
\newtheorem{definition}[subsection]{Definition}
\newtheorem{example}[subsection]{Example}
\newtheorem{exercise}[subsection]{Exercise}
\newtheorem{situation}[subsection]{Situation}

\theoremstyle{remark}
\newtheorem{remark}[subsection]{Remark}
\newtheorem{remarks}[subsection]{Remarks}

\numberwithin{equation}{subsection}

% Macros
%
\def\lim{\mathop{\mathrm{lim}}\nolimits}
\def\colim{\mathop{\mathrm{colim}}\nolimits}
\def\Spec{\mathop{\mathrm{Spec}}}
\def\Hom{\mathop{\mathrm{Hom}}\nolimits}
\def\Ext{\mathop{\mathrm{Ext}}\nolimits}
\def\SheafHom{\mathop{\mathcal{H}\!\mathit{om}}\nolimits}
\def\SheafExt{\mathop{\mathcal{E}\!\mathit{xt}}\nolimits}
\def\Sch{\mathit{Sch}}
\def\Mor{\mathop{\mathrm{Mor}}\nolimits}
\def\Ob{\mathop{\mathrm{Ob}}\nolimits}
\def\Sh{\mathop{\mathit{Sh}}\nolimits}
\def\NL{\mathop{N\!L}\nolimits}
\def\CH{\mathop{\mathrm{CH}}\nolimits}
\def\proetale{{pro\text{-}\acute{e}tale}}
\def\etale{{\acute{e}tale}}
\def\QCoh{\mathit{QCoh}}
\def\Ker{\mathop{\mathrm{Ker}}}
\def\Im{\mathop{\mathrm{Im}}}
\def\Coker{\mathop{\mathrm{Coker}}}
\def\Coim{\mathop{\mathrm{Coim}}}

% Boxtimes
%
\DeclareMathSymbol{\boxtimes}{\mathbin}{AMSa}{"02}

%
% Macros for moduli stacks/spaces
%
\def\QCohstack{\mathcal{QC}\!\mathit{oh}}
\def\Cohstack{\mathcal{C}\!\mathit{oh}}
\def\Spacesstack{\mathcal{S}\!\mathit{paces}}
\def\Quotfunctor{\mathrm{Quot}}
\def\Hilbfunctor{\mathrm{Hilb}}
\def\Curvesstack{\mathcal{C}\!\mathit{urves}}
\def\Polarizedstack{\mathcal{P}\!\mathit{olarized}}
\def\Complexesstack{\mathcal{C}\!\mathit{omplexes}}
% \Pic is the operator that assigns to X its picard group, usage \Pic(X)
% \Picardstack_{X/B} denotes the Picard stack of X over B
% \Picardfunctor_{X/B} denotes the Picard functor of X over B
\def\Pic{\mathop{\mathrm{Pic}}\nolimits}
\def\Picardstack{\mathcal{P}\!\mathit{ic}}
\def\Picardfunctor{\mathrm{Pic}}
\def\Deformationcategory{\mathcal{D}\!\mathit{ef}}
\usepackage{cleveref}
\crefname{originallemma}{original lemma}{original lemmas}
\Crefname{originallemma}{Original lemma}{Original lemmas}
\makeatletter
\newlabel{morphisms-lemma-base-change-proper}{{01W4}{source}{}{sourcepointer-1}{}}
\newlabel{morphisms-lemma-base-change-proper@cref}{{[originallemma][0][]01W4}{[1][0][]source}{}{}{}}
\makeatother
\crefname{originallemma}{original lemma}{original lemmas}
\Crefname{originallemma}{Original lemma}{Original lemmas}
\makeatletter
\newlabel{morphisms-lemma-closed-immersion-proper}{{01W5}{source}{}{sourcepointer-2}{}}
\newlabel{morphisms-lemma-closed-immersion-proper@cref}{{[originallemma][0][]01W5}{[1][0][]source}{}{}{}}
\makeatother
\crefname{originallemma}{original lemma}{original lemmas}
\Crefname{originallemma}{Original lemma}{Original lemmas}
\makeatletter
\newlabel{morphisms-lemma-composition-finite-type}{{01T3}{source}{}{sourcepointer-3}{}}
\newlabel{morphisms-lemma-composition-finite-type@cref}{{[originallemma][0][]01T3}{[1][0][]source}{}{}{}}
\makeatother
\crefname{originallemma}{original lemma}{original lemmas}
\Crefname{originallemma}{Original lemma}{Original lemmas}
\makeatletter
\newlabel{morphisms-lemma-composition-proper}{{01W3}{source}{}{sourcepointer-4}{}}
\newlabel{morphisms-lemma-composition-proper@cref}{{[originallemma][0][]01W3}{[1][0][]source}{}{}{}}
\makeatother
\crefname{originallemma}{original lemma}{original lemmas}
\Crefname{originallemma}{Original lemma}{Original lemmas}
\makeatletter
\newlabel{morphisms-lemma-image-proper-is-proper}{{03GN}{source}{}{sourcepointer-5}{}}
\newlabel{morphisms-lemma-image-proper-is-proper@cref}{{[originallemma][0][]03GN}{[1][0][]source}{}{}{}}
\makeatother
\crefname{originallemma}{original lemma}{original lemmas}
\Crefname{originallemma}{Original lemma}{Original lemmas}
\makeatletter
\newlabel{morphisms-lemma-image-universally-closed-separated}{{09MQ}{source}{}{sourcepointer-6}{}}
\newlabel{morphisms-lemma-image-universally-closed-separated@cref}{{[originallemma][0][]09MQ}{[1][0][]source}{}{}{}}
\makeatother
\crefname{originallemma}{original lemma}{original lemmas}
\Crefname{originallemma}{Original lemma}{Original lemmas}
\makeatletter
\newlabel{morphisms-lemma-immersion-locally-finite-type}{{01T5}{source}{}{sourcepointer-7}{}}
\newlabel{morphisms-lemma-immersion-locally-finite-type@cref}{{[originallemma][0][]01T5}{[1][0][]source}{}{}{}}
\makeatother
\crefname{originallemma}{original lemma}{original lemmas}
\Crefname{originallemma}{Original lemma}{Original lemmas}
\makeatletter
\newlabel{morphisms-lemma-scheme-theoretic-image-is-proper}{{0AH6}{source}{}{sourcepointer-8}{}}
\newlabel{morphisms-lemma-scheme-theoretic-image-is-proper@cref}{{[originallemma][0][]0AH6}{[1][0][]source}{}{}{}}
\makeatother
\crefname{originallemma}{original lemma}{original lemmas}
\Crefname{originallemma}{Original lemma}{Original lemmas}
\makeatletter
\newlabel{schemes-definition-inverse-image-closed-subscheme}{{01JV}{source}{}{sourcepointer-9}{}}
\newlabel{schemes-definition-inverse-image-closed-subscheme@cref}{{[originallemma][0][]01JV}{[1][0][]source}{}{}{}}
\makeatother
\title{Three permanence operations for closed subsets proper over a base: for f : X -> Y between S-schemes locally of finite type, images of S-proper closed subsets are S-proper closed when Y is separated over S or f is universally closed, and inverse images of S-proper closed subsets under proper f are S-proper closed.}
\author{The Stacks Project Authors}
\date{stacks-project:coherent.tex:local-snapshot-20261012; GNU FDL 1.2 or later}
\begin{document}
\maketitle
\noindent\textit{Editorial note.} Statement span is lines 7528-7541 of sources/stacks-coherent/pinned-original-source.tex; proof span is lines 7543-7588, the complete author proof (44 body lines).\par\medskip
\noindent\textit{Editorial note.} Navigation stubs: 9 cited-but-not-included labels are supplied as non-load-bearing pointers with their pinned tags where the snapshot tag index has them.\par\medskip
\noindent\textit{Editorial note.} extra\_packages is empty on purpose: the retained preamble already loads xy with options and a second declaration clashes in the assembler wrapper.\par\medskip
\noindent\textit{Editorial note.} Text is retained verbatim, including the author's own typos and placeholders.\par\medskip
\noindent\textit{Editorial note.} Pool stacks-coherent backfill wave 33c. precard\_gate RESULT=OK (probe L9999).\par\medskip
\noindent\textit{Allowed standard prerequisites (editorial).}\begin{itemize}
\item Explicit prerequisite boundary. See why-field; every pointer named here occurs verbatim in the retained statement/proof text.
\end{itemize}\medskip
\ifdefined\hypertarget\hypertarget{sourcepointer-1}{}\fi
\noindent\textit{Source reference.} Original lemma 01W4: sources/stacks-coherent/pinned-original-source.tex, source lines 7588--7588. Cited by the retained statement or proof but not part of the collected spans; supplied only as a navigation pointer. The argument does not depend on its content.\par\medskip
\ifdefined\hypertarget\hypertarget{sourcepointer-2}{}\fi
\noindent\textit{Source reference.} Original lemma 01W5: sources/stacks-coherent/pinned-original-source.tex, source lines 7588--7588. Cited by the retained statement or proof but not part of the collected spans; supplied only as a navigation pointer. The argument does not depend on its content.\par\medskip
\ifdefined\hypertarget\hypertarget{sourcepointer-3}{}\fi
\noindent\textit{Source reference.} Original lemma 01T3: sources/stacks-coherent/pinned-original-source.tex, source lines 7588--7588. Cited by the retained statement or proof but not part of the collected spans; supplied only as a navigation pointer. The argument does not depend on its content.\par\medskip
\ifdefined\hypertarget\hypertarget{sourcepointer-4}{}\fi
\noindent\textit{Source reference.} Original lemma 01W3: sources/stacks-coherent/pinned-original-source.tex, source lines 7588--7588. Cited by the retained statement or proof but not part of the collected spans; supplied only as a navigation pointer. The argument does not depend on its content.\par\medskip
\ifdefined\hypertarget\hypertarget{sourcepointer-5}{}\fi
\noindent\textit{Source reference.} Original lemma 03GN: sources/stacks-coherent/pinned-original-source.tex, source lines 7588--7588. Cited by the retained statement or proof but not part of the collected spans; supplied only as a navigation pointer. The argument does not depend on its content.\par\medskip
\ifdefined\hypertarget\hypertarget{sourcepointer-6}{}\fi
\noindent\textit{Source reference.} Original lemma 09MQ: sources/stacks-coherent/pinned-original-source.tex, source lines 7588--7588. Cited by the retained statement or proof but not part of the collected spans; supplied only as a navigation pointer. The argument does not depend on its content.\par\medskip
\ifdefined\hypertarget\hypertarget{sourcepointer-7}{}\fi
\noindent\textit{Source reference.} Original lemma 01T5: sources/stacks-coherent/pinned-original-source.tex, source lines 7588--7588. Cited by the retained statement or proof but not part of the collected spans; supplied only as a navigation pointer. The argument does not depend on its content.\par\medskip
\ifdefined\hypertarget\hypertarget{sourcepointer-8}{}\fi
\noindent\textit{Source reference.} Original lemma 0AH6: sources/stacks-coherent/pinned-original-source.tex, source lines 7588--7588. Cited by the retained statement or proof but not part of the collected spans; supplied only as a navigation pointer. The argument does not depend on its content.\par\medskip
\ifdefined\hypertarget\hypertarget{sourcepointer-9}{}\fi
\noindent\textit{Source reference.} Original lemma 01JV: sources/stacks-coherent/pinned-original-source.tex, source lines 7588--7588. Cited by the retained statement or proof but not part of the collected spans; supplied only as a navigation pointer. The argument does not depend on its content.\par\medskip
\begin{lemma}
\label{lemma-functoriality-closed-proper-over-base}
Let $S$ be a scheme. Let $f : X \to Y$ be a morphism of schemes which
are locally of finite type over $S$.
\begin{enumerate}
\item If $Y$ is separated over $S$ and $Z \subset X$ is a closed subset
proper over $S$, then $f(Z)$ is a closed subset of $Y$ proper over $S$.
\item If $f$ is universally closed and $Z \subset X$ is a
closed subset proper over $S$, then $f(Z)$ is a closed subset
of $Y$ proper over $S$.
\item If $f$ is proper and $Z \subset Y$ is a closed subset
proper over $S$, then $f^{-1}(Z)$ is a closed subset of $X$ proper over $S$.
\end{enumerate}
\end{lemma}

\begin{proof}
Proof of (1). Assume $Y$ is separated over $S$ and $Z \subset X$
is a closed subset proper over $S$. Endow $Z$ with the reduced induced
closed subscheme structure and apply
Morphisms, Lemma \ref{morphisms-lemma-scheme-theoretic-image-is-proper}
to $Z \to Y$ over $S$ to conclude.

\medskip\noindent
Proof of (2). Assume $f$ is universally closed and $Z \subset X$ is a
closed subset proper over $S$. Endow $Z$ and $Z' = f(Z)$ with their reduced
induced closed subscheme structures. We obtain an induced
morphism $Z \to Z'$.
Denote $Z'' = f^{-1}(Z')$ the scheme theoretic inverse image
(Schemes, Definition \ref{schemes-definition-inverse-image-closed-subscheme}).
Then $Z'' \to Z'$ is universally closed as a base change of $f$
(Morphisms, Lemma \ref{morphisms-lemma-base-change-proper}).
Hence $Z \to Z'$ is universally closed as a composition of
the closed immersion $Z \to Z''$ and $Z'' \to Z'$
(Morphisms, Lemmas
\ref{morphisms-lemma-closed-immersion-proper} and
\ref{morphisms-lemma-composition-proper}).
We conclude that $Z' \to S$ is separated by
Morphisms, Lemma \ref{morphisms-lemma-image-universally-closed-separated}.
Since $Z \to S$ is quasi-compact and $Z \to Z'$ is surjective
we see that $Z' \to S$ is quasi-compact.
Since $Z' \to S$ is the composition of $Z' \to Y$ and $Y \to S$
we see that $Z' \to S$ is locally of finite type
(Morphisms, Lemmas \ref{morphisms-lemma-immersion-locally-finite-type} and
\ref{morphisms-lemma-composition-finite-type}).
Finally, since $Z \to S$ is universally closed, we see that
the same thing is true for $Z' \to S$ by
Morphisms, Lemma \ref{morphisms-lemma-image-proper-is-proper}.
This finishes the proof.

\medskip\noindent
Proof of (3). Assume $f$ is proper and $Z \subset Y$ is a closed subset
proper over $S$. Endow $Z$ with the reduced induced closed subscheme
structure. Denote $Z' = f^{-1}(Z)$ the scheme theoretic inverse image
(Schemes, Definition \ref{schemes-definition-inverse-image-closed-subscheme}).
Then $Z' \to Z$ is proper as a base change of $f$
(Morphisms, Lemma \ref{morphisms-lemma-base-change-proper}).
Whence $Z' \to S$ is proper as the composition of $Z' \to Z$
and $Z \to S$
(Morphisms, Lemma \ref{morphisms-lemma-composition-proper}).
This finishes the proof.
\end{proof}

\end{document}
